\chapter{BEBERAPA TEORI SPEKTRAL}\label{spectrum}

Dalam bab ini, untuk pertama kalinya kita akan menganggap, kecuali dinyatakan lain, bahwa ruang
vektor dan aljabar berskalar kompleks, bukan real. Sebagian besar alasan untuk hal ini
 \index{konvensi!dalam bab~\ref{spectrum} skalar bersifat kompleks}%
adalah perlunya menggunakan \emph{teorema Liouville}~\ref{C073131} untuk membuktikan
proposisi~\ref{C073134}.

\section{Spektrum}

\begin{defn} Suatu elemen $a$ dari aljabar beridentitas $A$ disebut
 \index{invertibel!kiri!dalam aljabar}%
\df{invertibel kiri} jika terdapat elemen $a_l$ dalam $A$ (disebut \df{invers kiri}
dari~$a$) sedemikian sehingga $a_l a = \vc 1$, dan disebut
 \index{invertibel!kanan!dalam aljabar}%
\df{invertibel kanan} jika terdapat elemen $a_r$ dalam $A$ (disebut \df{invers
kanan} dari~$a$) sedemikian sehingga $aa_r = \vc 1$. Elemen tersebut disebut
 \index{invertibel!dalam aljabar}%
\df{invertibel} jika elemen itu sekaligus invertibel kiri dan invertibel kanan. Himpunan semua
elemen invertibel dari $A$ dinotasikan
 \index{invertibel@$\inv A$ (elemen invertibel dalam aljabar)}%
dengan~$\inv A$.
\end{defn}

Suatu elemen dari aljabar beridentitas memiliki paling banyak satu invers perkalian. Bahkan,
berlaku pernyataan yang lebih kuat berikut.

\begin{prop}\label{000521} Jika suatu elemen dari aljabar beridentitas memiliki invers kiri dan
invers kanan, maka kedua invers tersebut sama (sehingga elemen itu invertibel).
\end{prop}

Jika suatu elemen $a$ dari aljabar beridentitas invertibel, inversnya yang unik dinotasikan
dengan~$a^{-1}$.

\begin{prop} Jika $a$ elemen invertibel dari suatu aljabar beridentitas, maka $a^{-1}$ juga
 \index{invers!dari suatu invers}%
invertibel dan
   \[ \bigl(a^{-1}\bigr)^{-1} = a\,. \]
\end{prop}

\begin{prop} Jika $a$ dan $b$ elemen-elemen invertibel dari suatu aljabar beridentitas, maka hasil kali
 \index{invers!dari suatu hasil kali}%
$ab$ juga invertibel dan
  \[ (ab)^{-1} = b^{-1}a^{-1}\,. \]
\end{prop}

\begin{prop}\label{000524} Jika $a$ dan $b$ elemen-elemen invertibel dari suatu aljabar beridentitas,
 \index{invers!selisih antara}%
maka
  \[ a^{-1} - b^{-1} = a^{-1}(b - a)b^{-1}\,. \]
\end{prop}

\begin{prop} Misalkan $a$ dan $b$ elemen-elemen dari suatu aljabar beridentitas. Jika $ab$ dan
$ba$ keduanya invertibel, maka $a$ dan $b$ juga invertibel.
\end{prop}

\begin{proof}[\emph{Petunjuk untuk bukti}] Gunakan proposisi~\ref{000521}. \ns \end{proof}

\begin{cor} Misalkan $a$ elemen dari suatu aljabar-$*\,$ beridentitas. Maka $a$ invertibel jika dan hanya jika $a^*a$ dan $aa^*$
invertibel; dalam hal ini,
   \[ a^{-1} = \bigl(a^*a\bigr)^{-1}a^* = a^*\big(aa^*\bigr)^{-1}\,. \]
\end{cor}

\begin{prop}\label{000526} Misalkan $a$ dan $b$ elemen-elemen dari suatu aljabar beridentitas. Maka
$\vc 1 - ab$ invertibel jika dan hanya jika $\vc 1 - ba$ invertibel.
\end{prop}

\begin{proof} Jika $\vc 1 - ab$ invertibel, maka $\vc 1 + b(\vc 1 - ab)^{-1}a$ adalah
invers dari $\vc 1 - ba$.
\end{proof}

\begin{exam} Contoh paling sederhana dari aljabar kompleks ialah keluarga $\C$ dari bilangan
kompleks. Aljabar ini beridentitas sekaligus komutatif.
\end{exam}

\begin{exam}\label{000532} Jika $X$ suatu ruang topologis tak kosong, maka keluarga $\fml C(X)$
dari semua fungsi kontinu bernilai kompleks pada $X$ merupakan suatu
 \index{C@$\fml C(X)$!sebagai aljabar beridentitas}%
aljabar di bawah operasi penjumlahan, perkalian skalar, dan perkalian titik demi titik yang biasa.
Aljabar ini beridentitas sekaligus komutatif.
\end{exam}

\begin{exam} Jika $X$ suatu ruang Hausdorff kompak lokal yang tidak kompak, maka (di bawah
operasi titik demi titik) keluarga $\fml C_0(X)$ dari semua fungsi kontinu bernilai kompleks
pada $X$ yang lenyap di tak hingga merupakan aljabar komutatif.
Namun, keluarga ini \emph{bukan} suatu
 \index{C@$\fml C_0(X)$!sebagai aljabar tanpa identitas}%
aljabar beridentitas.
\end{exam}

\begin{exam}\label{000533} Keluarga
 \index{Ma@$\mathbf M_n = \M n\C$!matriks bilangan kompleks berukuran $n \times n$}%
$\mathbf M_n$ dari matriks bilangan kompleks berukuran $n \times n$ merupakan
 \index{Ma@$\mathbf M_n = \M n\C$!sebagai aljabar beridentitas}%
aljabar beridentitas di bawah operasi penjumlahan, perkalian skalar, dan perkalian matriks yang biasa.
Aljabar ini tidak komutatif jika $n > 1$.
\end{exam}

\begin{exam}\label{000533a} Misalkan $A$ suatu aljabar. Jadikan keluarga
 \index{Ma@$\M nA$!matriks elemen aljabar berukuran $n \times n$}%
$\M n{A}$ dari matriks elemen $A$ berukuran $n \times n$ sebagai suatu
 \index{Ma@$\M nA$!sebagai aljabar}%
aljabar dengan menggunakan aturan operasi matriks yang sama dengan aturan untuk~$\mathbf M_n$.
Dengan demikian, $\mathbf M_n$ hanyalah $\M n\C$. Aljabar $\M nA$ beridentitas jika dan hanya jika $A$ beridentitas.
\end{exam}

\begin{exam}\label{000534} Jika $V$ suatu ruang linear bernorma, maka
 \index{B@$\ofml B(V)$!sebagai aljabar beridentitas}%
$\ofml B(V)$ merupakan aljabar beridentitas. Jika $\dim V~>~1$, aljabar tersebut tidak komutatif.
\end{exam}

\begin{defn}\label{000535} Misalkan $a$ elemen dari aljabar beridentitas~$A$. Maka
 \index{spektrum}%
\df{spektrum} dari $a$, yang dinotasikan dengan
 \index{sigma@$\sigma_A(a)$ atau $\sigma(a)$ (spektrum dari~$a$)}%
$\sigma_A(a)$ atau cukup $\sigma(a)$, adalah himpunan semua bilangan kompleks $\lambda$ sedemikian sehingga
$a - \lambda\vc 1$ tidak invertibel.
\end{defn}

\begin{exam} Jika $z$ suatu elemen dari aljabar bilangan kompleks $\C$, maka
 \index{spektrum!dari bilangan kompleks}%
$\sigma(z) = \{z\}$.
\end{exam}

\begin{exam}\label{0005712} Misalkan $X$ suatu ruang Hausdorff kompak. Jika $f$ suatu elemen
 \index{spektrum!dari fungsi kontinu}%
dari aljabar $\fml C(X)$ yang terdiri atas fungsi-fungsi kontinu bernilai kompleks pada $X$, maka
spektrum $f$ adalah jangkauannya.
\end{exam}
